\section{Dissertation Structure}
\label{chap:introduction:structure}

The dissertation document is organized as follows. Chapter~\ref{chap:SOTA} presents the State of the Art covering aspects related to the \acs{5G} ecosystem, network slicing, \acs{NFV}, the \acs{5G} \acs{SBA} and exposure capabilities, network applications,  key telecom standardization initiatives, Management and Support Systems, as well as testbed implementations and testing \& validation procedures. Subsequently, Chapter~\ref{chap:proposal} outlines the dissertation's proposal, defining the testbed architecture, the capabilities it should provide for vertical experimentation, and a multi-vertical use case and validation methodology to evaluate the testbed. Chapter~\ref{chap:implementation} then describes the implementation of the proposed testbed and of the network application supporting the envisioned use case. Chapter~\ref{chap:results} presents the results of the validation, assessing the testbed and its exposure interfaces through the operation of the network application and its service delivery workflow. Finally, Chapter~\ref{chap:conclusions} summarizes the contributions of this dissertation and identifies directions for future work.
